% !TeX root = ../sections/03_solutions.tex
\subsection{1.1.B.6}

\begin{proof}
	分子を
	\[
	S_n=\sum_{j=1}^{n}j^2a_j
	\]
	とおく。
	
	まず，
	\[
	\frac{S_n}{n^3}\to \frac{\alpha}{3}
	\]
	であることを示す。
		仮定より
	\[
	a_j\to \alpha
	\]
	であるから，
	\[
	a_j=\alpha+(a_j-\alpha)
	\]
	と分解する。
	
	すると，
	\[
	S_n
	=
	\sum_{j=1}^{n}j^2a_j
	=
	\sum_{j=1}^{n}j^2\{\alpha+(a_j-\alpha)\}
	\]
	である。したがって，
	\[
	S_n
	=
	\alpha\sum_{j=1}^{n}j^2
	+
	\sum_{j=1}^{n}j^2(a_j-\alpha)
	\]
	となる。
	
	よって，
	\[
	\frac{S_n}{n^3}
	=
	\alpha\frac{1}{n^3}\sum_{j=1}^{n}j^2
	+
	\frac{1}{n^3}\sum_{j=1}^{n}j^2(a_j-\alpha)
	\]
	である。
		ここで，
	\[
	\sum_{j=1}^{n}j^2
	=
	\frac{n(n+1)(2n+1)}{6}
	\]
	であるから，
	\[
	\frac{1}{n^3}\sum_{j=1}^{n}j^2
	=
	\frac{n(n+1)(2n+1)}{6n^3}
	\]
	である。
	
	これを整理すると，
	\[
	\frac{n(n+1)(2n+1)}{6n^3}
	=
	\frac{2n^3+3n^2+n}{6n^3}
	\]
	であるから，
	\[
	\frac{1}{n^3}\sum_{j=1}^{n}j^2
	=
	\frac{1}{3}
	+
	\frac{1}{2n}
	+
	\frac{1}{6n^2}
	\]
	となる。
	
	したがって，
	\[
	\frac{1}{n^3}\sum_{j=1}^{n}j^2
	\to
	\frac{1}{3}
	\]
	である。
	
	よって，
	\[
	\alpha\frac{1}{n^3}\sum_{j=1}^{n}j^2
	\to
	\frac{\alpha}{3}
	\]
	である。
		次に，
	\[
	R_n=
	\frac{1}{n^3}\sum_{j=1}^{n}j^2(a_j-\alpha)
	\]
	とおく。
	
	$R_n\to 0$ を示す。
	
	任意に $\epsilon>0$ を取る。
	
	$a_j\to\alpha$ であるから，ある自然数 $N_1$ が存在して，
	$j\geq N_1$ ならば
	\[
	|a_j-\alpha|<\frac{\epsilon}{2}
	\]
	が成り立つ。
	
	有限個の和
	\[
	C=
	\sum_{j=1}^{N_1-1}j^2|a_j-\alpha|
	\]
	をおく。これは有限個の実数の和なので，有限の値である。
		$n\geq N_1$ のとき，三角不等式より
	\[
	|R_n|
	=
	\left|
	\frac{1}{n^3}\sum_{j=1}^{n}j^2(a_j-\alpha)
	\right|
	\leq
	\frac{1}{n^3}\sum_{j=1}^{n}j^2|a_j-\alpha|
	\]
	である。
	
	ここで和を分けると，
	\[
	\sum_{j=1}^{n}j^2|a_j-\alpha|
	=
	\sum_{j=1}^{N_1-1}j^2|a_j-\alpha|
	+
	\sum_{j=N_1}^{n}j^2|a_j-\alpha|
	\]
	である。
	
	したがって，
	\[
	|R_n|
	\leq
	\frac{C}{n^3}
	+
	\frac{1}{n^3}
	\sum_{j=N_1}^{n}j^2|a_j-\alpha|
	\]
	となる。
		$j\geq N_1$ ならば
	\[
	|a_j-\alpha|<\frac{\epsilon}{2}
	\]
	であるから，
	\[
	\sum_{j=N_1}^{n}j^2|a_j-\alpha|
	<
	\frac{\epsilon}{2}\sum_{j=N_1}^{n}j^2
	\leq
	\frac{\epsilon}{2}\sum_{j=1}^{n}j^2
	\]
	である。
	
	また，
	\[
	\sum_{j=1}^{n}j^2\leq n^3
	\]
	である。実際，$1\leq j\leq n$ ならば $j^2\leq n^2$ なので，
	$n$ 個足して高々 $n^3$ である。
	
	よって，
	\[
	\frac{1}{n^3}
	\sum_{j=N_1}^{n}j^2|a_j-\alpha|
	\leq
	\frac{\epsilon}{2}
	\]
	となる。
	
	したがって，$n\geq N_1$ のとき
	\[
	|R_n|
	\leq
	\frac{C}{n^3}
	+
	\frac{\epsilon}{2}
	\]
	である。
		さらに，$C/n^3\to0$ であるから，ある自然数 $N_2$ が存在して，
	$n\geq N_2$ ならば
	\[
	\frac{C}{n^3}<\frac{\epsilon}{2}
	\]
	が成り立つ。
	
	ここで
	\[
	N=\max\{N_1,N_2\}
	\]
	とおく。
	
	$n\geq N$ ならば
	\[
	|R_n|
	\leq
	\frac{C}{n^3}
	+
	\frac{\epsilon}{2}
	<
	\frac{\epsilon}{2}
	+
	\frac{\epsilon}{2}
	=
	\epsilon
	\]
	である。
	
	したがって
	\[
	R_n\to0
	\]
	である。
		以上より，
	\[
	\frac{S_n}{n^3}
	=
	\alpha\frac{1}{n^3}\sum_{j=1}^{n}j^2
	+
	R_n
	\]
	であり，
	\[
	\alpha\frac{1}{n^3}\sum_{j=1}^{n}j^2
	\to
	\frac{\alpha}{3},
	\qquad
	R_n\to0
	\]
	である。
	
	したがって，
	\[
	\frac{S_n}{n^3}\to\frac{\alpha}{3}
	\]
	である。
		いま求めたい極限は
	\[
	\lim_{n\to\infty}\frac{S_n}{n^k}
	\]
	である。
	
	すでに
	\[
	\frac{S_n}{n^3}\to\frac{\alpha}{3}
	\]
	が分かっている。
	
	また，仮定より $\alpha\neq0$ なので，
	\[
	\frac{\alpha}{3}\neq0
	\]
	である。
		まず $k=3$ のとき，
	\[
	\frac{S_n}{n^k}
	=
	\frac{S_n}{n^3}
	\to
	\frac{\alpha}{3}
	\]
	である。
	
	これは $0$ でない実数である。
		次に $k>3$ とする。
	
	このとき
	\[
	\frac{S_n}{n^k}
	=
	\frac{S_n}{n^3}\cdot\frac{1}{n^{k-3}}
	\]
	である。
	
	ここで
	\[
	\frac{S_n}{n^3}\to\frac{\alpha}{3}
	\]
	であり，
	\[
	\frac{1}{n^{k-3}}\to0
	\]
	である。
	
	したがって，
	\[
	\frac{S_n}{n^k}\to0
	\]
	となる。
	
	これは $0$ でない実数への収束ではない。
		最後に $k<3$ とする。
	
	このとき $3-k>0$ であり，
	\[
	\frac{S_n}{n^k}
	=
	\frac{S_n}{n^3}\cdot n^{3-k}
	\]
	である。
	
	ここで
	\[
	\frac{S_n}{n^3}\to\frac{\alpha}{3}
	\]
	であり，$\alpha\neq0$ なので，十分大きな $n$ に対して
	\[
	\left|\frac{S_n}{n^3}\right|
	\geq
	\frac{|\alpha|}{6}
	\]
	が成り立つ。
	
	一方，
	\[
	n^{3-k}\to\infty
	\]
	である。
	
	したがって，十分大きな $n$ に対して
	\[
	\left|\frac{S_n}{n^k}\right|
	=
	\left|\frac{S_n}{n^3}\right|n^{3-k}
	\geq
	\frac{|\alpha|}{6}n^{3-k}
	\]
	となる。
	
	右辺は $n\to\infty$ で正の無限大に発散するので，
	\[
	\frac{S_n}{n^k}
	\]
	は有限な実数には収束しない。
		以上より，
	\[
	\frac{S_n}{n^k}
	\]
	が $0$ でない実数に収束するためには，
	\[
	k=3
	\]
	でなければならない。
	
	そのときの極限値は
	\[
	A=\frac{\alpha}{3}
	\]
	である。
	
	したがって，
	\[
	\boxed{k=3,\qquad A=\frac{\alpha}{3}}
	\]
	である。
\end{proof}
\begin{lemma}
	任意の $n\in\mathbb{N}$ に対して，
	\[
	\sum_{j=1}^{n}j^2\leq n^3
	\]
	が成り立つ。
\end{lemma}

\begin{proof}
	$1\leq j\leq n$ とする。
	このとき
	\[
	j\leq n
	\]
	であるから，両辺は非負なので二乗して
	\[
	j^2\leq n^2
	\]
	を得る。
	
	したがって，
	\[
	1^2\leq n^2,\quad
	2^2\leq n^2,\quad
	\ldots,\quad
	n^2\leq n^2
	\]
	である。
	
	これらを足し合わせると，
	\[
	\sum_{j=1}^{n}j^2
	\leq
	\sum_{j=1}^{n}n^2
	\]
	となる。
	
	右辺は $n^2$ を $n$ 個足したものなので，
	\[
	\sum_{j=1}^{n}n^2=nn^2=n^3
	\]
	である。
	
	よって
	\[
	\sum_{j=1}^{n}j^2\leq n^3
	\]
	が成り立つ。
\end{proof}